%%%%%%%%%%%%%%%%%%%%%%%%%%%%%%%%%%%%%%%%%%%%%%%%%%%%%%%%%%%%%%%%%%%%%%
% Lớp 11 | Quan hệ vuông góc trong không gian | Hai mặt phẳng vuông góc
% Dạng 1H8?4-6: Bài toán cho trước góc giữa 𝑑và (𝑃)
% Số câu: 1  (sinh tự động bởi tach_cau.py; từ nay thầy sửa trực tiếp file này)
%%%%%%%%%%%%%%%%%%%%%%%%%%%%%%%%%%%%%%%%%%%%%%%%%%%%%%%%%%%%%%%%%%%%%%

%%%%%%%%%%%%%%%% TRẮC NGHIỆM 4 PHƯƠNG ÁN %%%%%%%%%%%%%%%%

% nguồn: Ok_TracNghiemLop11_Lan1
\begin{ex}%[1H8H4-6]%[ma:00143f59]
	Cho hình chóp $S.ABCD$  có đáy $ABCD$ là hình chữ nhật với $AB=a$, $AD=\sqrt{2}a$, $SA=3a$ và $SA\perp(ABCD)$. Góc giữa đường thẳng $SC$ và mặt phẳng $(ABCD)$ bằng
	 \choice
		 {\True $60^{\circ}$}
		{$120^{\circ}$} 
		{$30^{\circ}$} 
		 {$90^{\circ}$}
	 \loigiai{
	 \immini{
		Hình chiếu của $SC$  lên mặt phẳng $(ABCD)$ là $AC$ nên $\left(SC;(ABCD)\right)=\widehat{SCA}$.\\
		Ta có  $AC=\sqrt{AB^2+BC^2} =a\sqrt{3}$. \\
		$  \Rightarrow\tan\widehat{SAC}=\dfrac{SA}{AC} =\dfrac{3a}{a\sqrt{3}}=\sqrt{3}\Rightarrow\widehat{SCA}=60^{\circ} $.
		 }{
		\begin{tikzpicture}[scale=0.7,font=\footnotesize,line join=round,line cap=round,>=stealth]
				\path 
					 (0,0) coordinate (A)
					 (-1.3,-1.6) coordinate (B)
					(2.5,-1.6) coordinate (C) 
					($(A)+(C)-(B)$)  coordinate (D)
					($(A)+(0,3)$) coordinate (S); 
				\draw  (S)--(B)--(C)--(D)--cycle (S)--(C);
				 \draw[dashed] (S)--(A)--(D)
					 (C)--(A)--(B);
				\foreach \p/\q in {S/90,A/-90,B/-90,C/-90,D/0} 
				\fill[black]  (\p) circle (1.0pt) node[shift={(\q:2.5mm)}]{$\p$};
				\draw pic[draw,angle radius=0.2cm]  {right angle = S--A--D};
			\end{tikzpicture}
	}} 
\end{ex}
